
\documentclass{article}
\usepackage{colortbl}
\usepackage{makecell}
\usepackage{multirow}
\usepackage{supertabular}

\begin{document}

\newcounter{utterance}

\centering \large Interaction Transcript for game `cladder', experiment `full\_v1.5\_default', episode 9977 with selfrefine{-}llama{-}3.1{-}8b{-}instruct.
\vspace{24pt}

{ \footnotesize  \setcounter{utterance}{1}
\setlength{\tabcolsep}{0pt}
\begin{supertabular}{c@{$\;$}|p{.15\linewidth}@{}p{.15\linewidth}p{.15\linewidth}p{.15\linewidth}p{.15\linewidth}p{.15\linewidth}}
    \# & \multicolumn{2}{c}{Player} && \multicolumn{2}{c}{Game Master} \\
    \hline

    \theutterance \stepcounter{utterance}  
    & & & \multicolumn{4}{p{0.6\linewidth}}{
        \cellcolor[rgb]{0.9,0.9,0.9}{
            \makecell[{{p{\linewidth}}}]{
                \texttt{\tiny{[P1$\langle$GM]}}
                \texttt{You are an expert in causal inference. The following question is not a typical commonsense query, but rather a meticulously designed question created by a professor specializing in causal inference, intended to assess the students' mastery of the course content.} \\
\texttt{Imagine a self{-}contained, hypothetical world with only the following conditions, and without any unmentioned factors or causal relationships: Hwax has a direct effect on pexu and rukz. Pexu has a direct effect on rukz. We know that hwax causes pexu. hwax or pexu causes rukz. We observed an individual is not hwax. Would an individual is rukz if pexu instead of not pexu?} \\
\texttt{Start your answer with 'Yes' or 'No', followed by additional reasoning or evidence to support your explanation.} \\
            }
        }
    }
     \\ \\

    \theutterance \stepcounter{utterance}  
    & & \multicolumn{4}{p{0.6\linewidth}}{
        \cellcolor[rgb]{0.95,0.95,0.95}{
            \makecell[{{p{\linewidth}}}]{
                \texttt{\tiny{[P1$\leftrightarrow$P]}}
                \texttt{{[}DRAFT{]}} \\
\texttt{No.} \\
\\ 
\texttt{The given conditions state that hwax causes pexu, and both hwax and pexu can cause rukz. However, the question asks us to consider a scenario where the individual is not hwax, but instead, pexu is present. } \\
\\ 
\texttt{To determine the causal effect of pexu on rukz in the absence of hwax, we need to consider the causal relationships between these variables. Since pexu has a direct effect on rukz, and we know that pexu is present, it is likely that pexu causes rukz.} \\
\\ 
\texttt{However, the key point here is that the individual is not hwax. The condition states that hwax has a direct effect on pexu and rukz. This implies that hwax and pexu are not independent causes of rukz. In other words, the effect of pexu on rukz is mediated by hwax.} \\
\\ 
\texttt{Given that the individual is not hwax, we cannot assume that pexu would cause rukz independently of hwax. The presence of pexu may still be influenced by the absence of hwax, and the causal effect of pexu on rukz may be altered as a result.} \\
\\ 
\texttt{Therefore, based on the given conditions, we cannot conclude that an individual would be rukz if pexu instead of not pexu, in the absence of hwax.} \\
            }
        }
    }
    & \\ \\

    \theutterance \stepcounter{utterance}  
    & & \multicolumn{4}{p{0.6\linewidth}}{
        \cellcolor[rgb]{0.95,0.95,0.95}{
            \makecell[{{p{\linewidth}}}]{
                \texttt{\tiny{[P1$\leftrightarrow$P]}}
                \texttt{{[}SELF{-}FEEDBACK{]}} \\
\texttt{FORMAT: ok} \\
\texttt{RULES: ok} \\
\texttt{STRATEGY: problem} \\
\texttt{REASON: The answer does not directly address the question, instead providing a detailed explanation of the causal relationships between the variables, which is not necessary to answer the question.} \\
            }
        }
    }
    & \\ \\

    \theutterance \stepcounter{utterance}  
    & \multicolumn{4}{p{0.6\linewidth}}{
        \cellcolor[rgb]{0.561,0.737,0.561}{
            \makecell[{{p{\linewidth}}}]{
                \texttt{\tiny{[P1$\rangle$GM]}}
                \texttt{No.} \\
\\ 
\texttt{The given conditions state that hwax causes pexu, and both hwax and pexu can cause rukz. However, the question asks us to consider a scenario where the individual is not hwax, but instead, pexu is present. } \\
\\ 
\texttt{Since pexu has a direct effect on rukz, and we know that pexu is present, it is likely that pexu causes rukz. However, the key point here is that the individual is not hwax. } \\
\\ 
\texttt{Given that the individual is not hwax, we cannot assume that pexu would cause rukz independently of hwax. The presence of pexu may still be influenced by the absence of hwax, and the causal effect of pexu on rukz may be altered as a result.} \\
\\ 
\texttt{Therefore, based on the given conditions, we cannot conclude that an individual would be rukz if pexu instead of not pexu, in the absence of hwax.} \\
            }
        }
    }
    & & \\ \\

    \theutterance \stepcounter{utterance}  
    & & & \multicolumn{2}{p{0.3\linewidth}}{
        \cellcolor[rgb]{0.9,0.9,0.9}{
            \makecell[{{p{\linewidth}}}]{
                \texttt{\tiny{[GM$|$GM]}}
                \texttt{No} \\
            }
        }
    }
    & & \\ \\

    \theutterance \stepcounter{utterance}  
    & & & \multicolumn{2}{p{0.3\linewidth}}{
        \cellcolor[rgb]{0.9,0.9,0.9}{
            \makecell[{{p{\linewidth}}}]{
                \texttt{\tiny{[GM$|$GM]}}
                \texttt{yes} \\
            }
        }
    }
    & & \\ \\

    \theutterance \stepcounter{utterance}  
    & & & \multicolumn{2}{p{0.3\linewidth}}{
        \cellcolor[rgb]{0.9,0.9,0.9}{
            \makecell[{{p{\linewidth}}}]{
                \texttt{\tiny{[GM$|$GM]}}
                \texttt{game\_result = LOSE} \\
            }
        }
    }
    & & \\ \\

\end{supertabular}
}

\end{document}
